\section{A new unconditional local reduction from 238 to 237 states}
\label{sec:local237}

The last state saving uses a further coordinated helper change.  Its proof is
particularly short: three replacement blocks preserve the tape at every
intermediate step, and all their guards follow from an inductive predecessor
invariant.  This reduction requires no new assumption about valid register
frames.  It therefore proves halting equivalence from the inherited start state
on \emph{every} initial binary tape.

Throughout this section, states retain their inherited numeric labels.
Let $M$ be the 238-state input machine and $N$ the new 237-state machine.
Both start in state 95 and halt on entering state 254.  A row written
$wDq$ writes $w$, moves in direction $D$, and enters state $q$.
The 238-state input has SHA256
\begin{quote}\small\ttfamily
069942be27db9736875ea042320e51ac3b1\\
eac7116cbe01cdce520fa6056f78c.
\end{quote}
The 237-state output has inherited-label SHA256
\begin{quote}\small\ttfamily
a3c4407f30e2e2139af87c6a91948c5181\\
34855841a28eb85b1d81bcfd026c9a.
\end{quote}

\subsection{Exact construction}

Delete state 235 and redirect its incoming edges to state 98.  Change one row
of state 98 and one row of the retained helper state 240.  The complete list
of changes to retained rows is
\begin{center}
\begin{tabular}{@{}lll@{}}
\toprule
Row & $M$ & $N$\\
\midrule
$98/0$ & $0R113$ & $0L240$\\
$139/1$ & $1R235$ & $1R98$\\
$156/1$ & $1R235$ & $1R98$\\
$240/1$ & $1R99$ & $1L87$\\
\bottomrule
\end{tabular}
\end{center}
The deleted rows are $235/0=0R231$ and $235/1=1R16$.
The retained row $98/1=1R16$ is unchanged.  Thus all 470 other retained rows
are unchanged, and the output has exactly 474 transition rows on 237 working
states.  The separate halt state is not included in the working-state count.
The source program, alphabet, tape convention, initial state, and initializer
are inherited without an additional input word or an uncharged control state.

Define the map on control states by
\[
 \phi(q)=\begin{cases}98,&q=235,\\q,&q\ne235.\end{cases}
\]
It fixes state 95 and halt state 254.  For a configuration
$C=(q,h,T)$, where $T:\mathbb Z\to\{0,1\}$, write
$F(C)=(\phi(q),h,T)$.

\subsection{The predecessor invariant}

For an old configuration $(q,h,T)$, let $I(q,h,T)$ be the conjunction of
the following implications:
\begin{align*}
 q=98 &\ \Longrightarrow\ T(h-1)=0,\\
 q=139 &\ \Longrightarrow\ T(h-1)=1,\\
 q=156 &\ \Longrightarrow\ T(h-1)=1,\\
 q=235 &\ \Longrightarrow\ T(h-2)=T(h-1)=1,\\
 q=240 &\ \Longrightarrow\ T(h-1)=0.
\end{align*}
This invariant concerns only the run of $M$.  It need not hold at new
intermediate configurations of $N$, where the helper intentionally serves
another role.

The incoming-edge inventory of the exact old table is
\begin{center}
\begin{tabular}{@{}lll@{}}
\toprule
Target & Incoming rows & Immediate consequence\\
\midrule
98 & $86/0$, $187/0$ & Write 0 and move right\\
139 & $117/1$ & Write 1 and move right\\
156 & $195/1$ & Write 1 and move right\\
235 & $139/1$, $156/1$ & Write 1 and move right\\
240 & $135/0$, $242/0$ & Write 0 and move right\\
\bottomrule
\end{tabular}
\end{center}
The first, second, third, and fifth implications follow directly from these
writes and moves.  On entry to 235, the new immediate-left bit is the 1 just
written.  The old invariant at predecessor 139 or 156 supplies the new
two-left bit, also 1.  Hence every old transition preserves $I$.
State 95 is none of the five constrained states, so $I(95,h,T)$ holds for
every $h$ and every binary tape $T$.  Induction on the old step count proves
$I$ throughout every such old run.

The accompanying checker independently verifies this implication over all
nine-cell Boolean windows admitted by $I$, with arbitrary exterior tape.
It closes 120,448 state/window cases.  The nine-cell window is more than the
proof needs; every read and every invariant query lies within the represented
window.  This is a finite check of an inductive statement, not a sample of
possible executions.

A second independent checker uses entailment of finite partial tape
assignments.  It verifies invariant preservation for all 476 old transition
rows directly, without enumerating the first checker's nine-cell windows.
It also proves that the invariant and read bit entail every replacement
block guard, and checks the three open-tape blocks and all 473 ordinary
row identities.  Its complete result is retained as
\file{evidence/local237_independent_review.json}.

\subsection{Three open-tape block identities}

All nontrivial old steps are accounted for by the following table.  Coordinates
are relative to the initial head position, which is set to 0 only for notation.
The final state in each displayed path is the synchronization endpoint.
\begin{center}
\begin{tabular}{@{}llll@{}}
\toprule
Old row & Required left bits & New control path & Steps\\
\midrule
$98/0$ & $T(-1)=0$
 & $98\to240\to176\to113$ & 3\\
$235/0$ & $T(-2)=T(-1)=1$
 & $98\to240\to87\to164\to118\to231$ & 5\\
$240/1$ & $T(-1)=0$
 & $240\to87\to190\to99$ & 3\\
\bottomrule
\end{tabular}
\end{center}
The guards are exactly the relevant consequences of $I$.
For complete transparency, the read configurations of the three new blocks
are, respectively,
\begin{align*}
 &(98,0,0),\ (240,-1,0),\ (176,0,0),\\
 &(98,0,0),\ (240,-1,1),\ (87,-2,1),\
   (164,-1,1),\ (118,0,0),\\
 &(240,0,1),\ (87,-1,0),\ (190,0,1).
\end{align*}
Each triple lists state, relative head position, and bit read.  Every row
on these paths writes back the bit it reads.  Thus the tape is unchanged
even between synchronization points.  The final head position is $+1$ in
each case, and the final states are 113, 231, and 99, exactly as required
by the respective old steps.  No cell outside $[-2,0]$ is read or written.

The checker proves these identities on symbolic open tapes: each unbound
cell is an independent Boolean variable, and it rejects any attempted branch
on an unbound cell.  The displayed guards resolve every branch.  The final
whole-tape override map agrees exactly with that of the old step.

The remaining 473 old transition rows commute in one step through $F$.
This includes $235/1$, whose transition is identical to retained $98/1$.
The two incoming-edge redirections simply implement $\phi$ on the target.
Therefore every old reachable nonhalting configuration $C$ satisfies
\[
 N^{k(C)}(F(C))=F(M(C)),
 \qquad
 k(C)=
 \begin{cases}
 3,&(q,T(h))=(98,0),\\
 5,&(q,T(h))=(235,0),\\
 3,&(q,T(h))=(240,1),\\
 1,&\text{otherwise}.
 \end{cases}
\]
Here $M(C)$ denotes one transition, and $N^k$ denotes $k$ successive
transitions.  None of the exceptional blocks has an interior halt.

\subsection{Halting equivalence and exact clock relation}

\paragraph{Theorem.}
For every binary tape $T$ and every head position $h$,
\[
 \operatorname{Halt}(M,95,h,T)
 \quad\Longleftrightarrow\quad
 \operatorname{Halt}(N,95,h,T).
\]
The proof uses only the exact finite tables and induction.  With tapes
represented as functions on the integers, this theorem can in particular
be formalized in ZFC.

\paragraph{Proof.}
Let $C_t$ be the old run while it is defined, and put $\tau(0)=0$ and
$\tau(t+1)=\tau(t)+k(C_t)$.  Induction using the preceding block identities
gives $N^{\tau(t)}(F(C_0))=F(C_t)$.  An old halt is thus a new halt.
If the old run is infinite, all its synchronization states are nonhalting,
each intervening block has no interior halt, and $\tau(t)\ge t$.
The new blocks consequently cover every finite prefix of an infinite new
run.  The new machine cannot halt.  This proves the reverse implication
as well.  The argument uses only positive blocks of length at most five.
\hfill$\square$

More precisely, if $a_t$, $b_t$, and $c_t$ count visits to old rows
$98/0$, $235/0$, and $240/1$ among the first $t$ old transitions, then
\[
 \tau(t)=t+2a_t+4b_t+2c_t,
 \qquad t\le\tau(t)\le5t.
\]
This is an exact identity for each prefix, not an asymptotic estimate.
The supplied startup receipt instantiates it for the complete prescribed
238-state blank initialization.  None of the three exceptional rows is
visited in that prefix, so the new startup clock is unchanged at 409,835
steps.  The checker compares all 409,836 synchronization checkpoints,
including the initial and final configurations, and independently checks
the final tape with sparse and dense emulators.  The endpoint has inherited
state 96, head $-5$, marker $-6$, header 176, 351 one-cells, logical registers
$(3,0,0,0,3,0,0,0,0)$, and 167 spare unary blocks of value 1.  All other tape
cells are zero.

Applying the theorem to the blank tape transfers the established 238-state
calibration immediately:
\[
 \mathrm{ZFC}\ \vdash\
 \neg\operatorname{Halt}(N,\mathrm{blank})
 \ \longleftrightarrow\ \operatorname{Con}(\mathrm{ZFC}).
\]
This last displayed conclusion inherits the preceding source-calibration
and 238-state correctness premises.  The new local argument itself has no
frame-sensitive guard.  It supplies neither an unconditional proof of
nonhalting nor a proof-assistant kernel theorem for the final machine.

\subsection{What the new reduction adds}

The 238-state combination changes the decoder contexts enough to expose a
further sharing opportunity.  Old states 98 and 235 have the same read-1
transition.  On read 0 their left contexts are disjoint: 0 for state 98 and
11 for state 235.  The revised row of helper 240 distinguishes these roles
while preserving its original behavior in its own left-0 contexts.
The machine stores the distinction in already-present tape bits and pays
for it by taking a few extra steps.  No marker bit, extra tape alphabet
symbol, hidden input, or added state is introduced.

Finite instruction profiles guided discovery only.  The accepted saving
rests on the predecessor invariant and the three open-tape identities above;
their scope is stronger than the discovery profile.  A separate complete
unbounded instruction-runtime replay of the exact 237-state table supplies
an additional validation route: all 8,386 guarded cases close, using
2,261,369 symbolic rules, with minimum instruction duration 13 and all 474
rows exercised.  The bounded search that found this edit
does not establish a state minimum or rule out other encodings.
